\documentclass[11pt]{article}
\usepackage[margin=1in]{geometry}
\usepackage{amsmath,amssymb}
\usepackage{booktabs}
\usepackage{hyperref}

\title{PGD Attack and Defense: Implementation and Analysis}
\author{Assignment 3 - Part 2}
\date{}

\begin{document}
\maketitle

\section{Question 1: PGD Implementation and Findings}

\subsection{Implementation}

\textbf{PGD Attack (\texttt{pgd\_attack} in \texttt{adversarial\_attack.py}):}
The PGD attack iteratively perturbs the input and projects back onto the $\epsilon$-ball:
\begin{verbatim}
def pgd_attack(model, data, target, criterion, args):
    x_orig = data.detach()
    x_adv = x_orig.clone()
    for _ in range(num_iter):
        x_adv.requires_grad_(True)
        model.zero_grad()
        loss = criterion(model(x_adv), target)
        loss.backward()
        with torch.no_grad():
            x_adv = x_adv + alpha * x_adv.grad.sign()
            perturbation = torch.clamp(x_adv - x_orig, -epsilon, epsilon)
            x_adv = torch.clamp(x_orig + perturbation, 0, 1)
        x_adv = x_adv.detach()
    return x_adv
\end{verbatim}

\textbf{PGD in \texttt{test\_attack}:}
\begin{verbatim}
elif attack_function == PGD:
    perturbed_data = pgd_attack(model, data, target, criterion, attack_args)
    output = model(perturbed_data)
\end{verbatim}

\textbf{PGD Defense in \texttt{utils.py}:}
\begin{verbatim}
elif defense_strategy == PGD and phase == 'train':
    adv_inputs = pgd_attack(model, inputs, labels, criterion, defense_args)
    inputs = torch.cat([inputs, adv_inputs], dim=0)
    labels = torch.cat([labels, labels], dim=0)
    optimizer.zero_grad()
    outputs = model(inputs)
    loss = criterion(outputs, labels)
    _, preds = torch.max(outputs, 1)
\end{verbatim}

\subsection{Findings}

\begin{table}[h]
\centering
\begin{tabular}{lccc}
\toprule
\textbf{Configuration} & \textbf{Clean Acc.} & \textbf{FGSM Acc.} & \textbf{PGD Acc.} \\
\midrule
No Defense (pretrained) & 92.96\% & 41.68\% & 21.40\% \\
FGSM Defense (5 epochs) & 89.88\% & 58.56\% & 43.68\% \\
\bottomrule
\end{tabular}
\caption{Configuration: $\epsilon_{\text{FGSM}}=0.1$, $\epsilon_{\text{PGD}}=0.03$, $\alpha_{\text{PGD}}=0.007$, 10 iterations.}
\end{table}

PGD is a stronger attack than FGSM: under PGD attack, the model retains only 21.40\% accuracy compared to 41.68\% under FGSM. This is because PGD iteratively refines the perturbation to find adversarial examples closer to the worst-case within the $\epsilon$-ball, whereas FGSM takes only a single step.

\section{Question 2: Adversarial Loss vs. Adding Adversarial Examples to Batch}

The two defense approaches are \textbf{different}, though they share the same goal of improving robustness.

\textbf{Adversarial Loss (FGSM-style):}
\[
\mathcal{L}_{\text{total}} = \mathcal{L}(x, y) + \alpha \cdot \mathcal{L}(x_{\text{adv}}, y)
\]
Both terms are combined before backpropagation. The parameter $\alpha$ controls the relative importance of adversarial robustness versus clean accuracy.

\textbf{Adding to Batch (PGD-style):}
The batch is augmented with adversarial examples, then standard loss is computed:
\[
\mathcal{L} = \frac{1}{2N}\sum_{i=1}^{N}\left[\mathcal{L}(x_i, y_i) + \mathcal{L}(x_{\text{adv},i}, y_i)\right]
\]

\textbf{When they align:} If $\alpha = 1$, both methods weight clean and adversarial examples equally, producing equivalent gradient contributions.

\textbf{When they diverge:}
\begin{enumerate}
    \item Adversarial loss allows flexible weighting via $\alpha$; batch augmentation treats both equally.
    \item Batch augmentation doubles the effective batch size, affecting batch normalization statistics and optimizer dynamics.
    \item In FGSM loss, adversarial examples share the computation graph with the gradient update. In batch augmentation, PGD examples are pre-computed and detached before the forward pass.
\end{enumerate}

\section{Question 3: FGSM vs. PGD Tradeoff}

\begin{table}[h]
\centering
\begin{tabular}{lcc}
\toprule
\textbf{Aspect} & \textbf{FGSM} & \textbf{PGD} \\
\midrule
Gradient computations & 1 & $k$ (number of iterations) \\
Attack strength & Weaker & Stronger \\
Training time & Fast & $\sim k\times$ slower \\
Robustness achieved & Moderate & Higher \\
\bottomrule
\end{tabular}
\end{table}

\textbf{FGSM Advantage:} Computational efficiency. FGSM requires only one forward and one backward pass, making it practical for large-scale adversarial training where computational resources are limited.

\textbf{PGD Advantage:} Stronger adversarial examples and better robustness. PGD iteratively refines the perturbation to find examples closer to the worst-case within the $\epsilon$-ball. Models trained against PGD attacks are more robust because they learn to defend against a stronger adversary.

\end{document}
